\section{实战落地：事件复盘与信任恢复}

\subsection{为什么“复盘”是 AI SRE 成熟的第二阶段}

事件处理完并不意味着工作结束。关键在于“复盘和信任恢复”：

\begin{itemize}
    \item 记录所有 Agent 交互、工具调用、自动建议和人工决策
    \item 分析为什么 AI 没能给出正确建议（数据缺失、prompt 注入、模型 hallucination）
    \item 补回缺失的文档、hooks 和监控信号
    \item 向团队公布经验教训，更新 runbook 和审批政策
\end{itemize}

\begin{warningbox}{没有复盘就没有信任}
运行中的团队对 AI 的信任来自连续的成功案例和透明的错误处理路径。如果每次 AI 出问题都只是“回滚并忘记”，团队就不会允许它承担更多权限。
\end{warningbox}

\subsection{事件复盘的可视化流程}

\begin{table}[H]
\centering
\begin{tabular}{p{0.2\textwidth} p{0.25\textwidth} p{0.28\textwidth}}
\toprule
步骤 & 产出 & 关键指标 \\
\midrule
收集证据 & 命令历史 + prompt + tool outputs & 覆盖率、日志完整性 \\
归因分析 & 根因链 + 证据矩阵 & 误差区间、关联度 \\
补救方案 & 新 hook / 新测试 / 调整权限 & 执行次数、Rollback 次数 \\
知识沉淀 & 更新 runbook、加标签、分享会议 & 文档完成度、审计验证 \\
\bottomrule
\end{tabular}
\caption{事件复盘的可视化流程}
\end{table}

\subsection{本章小结}

事件处理的最后一公里是复盘与信任恢复。记录原始决策链条、分析根因、修补监控，并把教训转化为可操作的 runbook，是 AI SRE 真正进入生产的前提。

